\section*{Hoofdstuk \ref{ch:b1:plate-kinematics} --- Plaatbeweging op een bol}

\begin{solution}{exo:b1:plate-kinematics:1}
(a) $\qty{0.210}{\degree\per\Myr} = 0.210\times\pi/180 = \qty{3.665e-3}{rad}$
per miljoen jaar, dat is \qty{3.665e-9}{rad.yr^{-1}}; met
$\qty{1}{yr} = \qty{3.156e7}{s}$, \qty{1.161e-16}{rad.s^{-1}}. (b) $360/0.210 =
\qty{1714}{\Myr}$, meer dan een derde van de leeftijd van de aarde. (c) $\omega R =
\num{3.665e-9}\times\qty{6.371e9}{mm} = \qty{23.35}{mm.yr^{-1}}$.
\end{solution}

\begin{solution}{exo:b1:plate-kinematics:2}
$v = \omega R\sin\Delta = 23.35\sin\Delta$: \qty{11.7}{mm.yr^{-1}} bij \ang{30},
\qty{20.2}{mm.yr^{-1}} bij \ang{60} en \qty{23.35}{mm.yr^{-1}} bij \ang{90}. De
spreiding is het snelst waar de grens het verst van de pool ligt, tot
\ang{90}: tussen Eurazië en Noord-Amerika, aan het zuidelijke einde van de grens, nabij de Azoren, \ang{78} van de pool.
\end{solution}

\begin{solution}{exo:b1:plate-kinematics:3}
De subchron duurde $1.072 - 0.988 = \qty{0.084}{\Myr}$, en een halve snelheid in
\unit{mm.yr^{-1}} is ook een snelheid in \unit{km\per\Myr}. (a) $50\times0.084 =
\qty{4.2}{km}$ bij \qty{50}{mm.yr^{-1}}; $10\times0.084 = \qty{0.84}{km}$ bij
\qty{10}{mm.yr^{-1}}. (b) De streep omvat de as:
$2\times10\times0.781 = \qty{15.6}{km}$.
\end{solution}

\begin{solution}{exo:b1:plate-kinematics:4}
De halve snelheden zijn $66/2.581 = \qty{25.6}{mm.yr^{-1}}$ naar het westen en
$64/2.581 = \qty{24.8}{mm.yr^{-1}}$ naar het oosten; de volle snelheid is
$130/2.581 = \qty{50.4}{mm.yr^{-1}}$. De
spreiding is bijna symmetrisch: de westelijke flank is ongeveer \qty{3}{\percent} sneller gegroeid.
\end{solution}

\begin{solution}{exo:b1:plate-kinematics:5}
$\cos\Delta = \sin\ang{61.8}\sin\ang{85} + \cos\ang{61.8}\cos\ang{85}
\cos(\ang{0} - \ang{139.6}) = 0.8813\times0.9962 + 0.4726\times0.08716\times
(-0.7615) = 0.87795 - 0.03137 = 0.84658$, dus $\Delta = \ang{32.16}$ en
$v = 23.35\sin\ang{32.16} = \qty{12.4}{mm.yr^{-1}}$. Dit ligt iets boven de gemeten \qtyrange{11}{12}{mm.yr^{-1}}, maar de \omterm{def:b1:plate-kinematics:euler-pole}{Eulervector} is één beste schatting aan 453 snelheden, waarvan de snelheden gemiddeld \qty{0.7}{mm.yr^{-1}} afwijken, en afzonderlijke snelheden nabij de trage Gakkelrug behoren tot de minst nauwkeurige.
\end{solution}

\begin{solution}{exo:b1:plate-kinematics:6}
(a) $\cos\Delta = \sin\ang{61.8}\sin\ang{39.44} + \cos\ang{61.8}\cos\ang{39.44}
\cos(\ang{169.43}) = 0.20111$: $\Delta = \ang{78.40}$. (b) $v = 23.35
\sin\ang{78.40} = \qty{22.87}{mm.yr^{-1}}$, de \qty{23}{mm.yr^{-1}} die net ten noorden van het knooppunt wordt gemeten. (c) Met de formule van de methode is de teller
$\sin\ang{169.43}\cos\ang{61.8} = 0.08668$ en de noemer
$\cos\ang{39.44}\sin\ang{61.8} - \sin\ang{39.44}\cos\ang{61.8}\cos\ang{169.43}
= 0.97573$: de pool wordt gezien onder $\beta = \ang{5.08}$, bijna recht naar het noorden, en Eurazië beweegt naar $\ang{95.1}$. (d) De transforms lopen evenwijdig aan de relatieve beweging: zij lopen ongeveer langs \ang{095}, vrijwel oost-west.
\end{solution}

\begin{solution}{exo:b1:plate-kinematics:7}
(a) $\vec\omega = \omega(\cos\varphi_p\cos\lambda_p, \cos\varphi_p\sin\lambda_p,
\sin\varphi_p)$: $\vec\omega_{\mathrm{EU\,NA}} = (-0.0756, 0.0643, 0.1851)$ en
$\vec\omega_{\mathrm{NB\,NA}} = (0.0333, 0.0282, 0.2289)$, in
\unit{\degree\per\Myr}. (b) $\vec\omega_{\mathrm{NB\,EU}} =
\vec\omega_{\mathrm{NB\,NA}} - \vec\omega_{\mathrm{EU\,NA}} = (0.1089, -0.0361,
0.0438)$, met lengte
\[
  \sqrt{0.011859 + 0.001303 + 0.001918} = \qty{0.1228}{\degree\per\Myr}.
\]
De \omterm{def:b1:plate-kinematics:euler-pole}{Eulerpool} ligt dus op de noordelijke \omterm{def:b1:spherical-astronomy:geographic}{breedtegraad}
$\arcsin(0.0438/0.1228) = \ang{20.9}$\,N en \omterm{def:b1:spherical-astronomy:geographic}{lengtegraad} $\arctan(-0.0361/0.1089) = \ang{-18.3}$, dat is \ang{18.3}\,W (de $x$-component is positief, dus er wordt geen \ang{180} bij opgeteld). (c) De grens tussen Nubië en Eurazië biedt weinig directe metingen: tussen de Azoren en Gibraltar alleen de \omterm{def:b1:spherical-astronomy:horizontal}{azimuts} van de Gloriabreuk. De twee spreidingsbewegingen over de Atlantische Oceaan worden met honderden snelheden gemeten, en hun verschil erft hun nauwkeurigheid.
\end{solution}

\begin{solution}{exo:b1:plate-kinematics:8}
(a) Om $\vec e_z$ over \ang{90}: $\cos\ang{90} = 0$, $\vec e_z\cdot\vec e_x = 0$
en $\vec e_z\times\vec e_x = \vec e_y$, dus $\vec e_x$ gaat naar $\vec e_y$. Dan om $\vec e_x$ over \ang{90}: $\vec e_x\cdot\vec e_y = 0$ en
$\vec e_x\times\vec e_y = \vec e_z$, dus $\vec e_y$ gaat naar $\vec e_z$: het punt eindigt op de noordpool. (b) Eerst om $\vec e_x$: $\vec e_x$ ligt op de as,
$\vec e_x\times\vec e_x = \vec 0$ en $(1 - \cos\ang{90})(\vec e_x\cdot\vec e_x)
\vec e_x = \vec e_x$, dus het beweegt niet; dan om $\vec e_z$ gaat het naar
$\vec e_y$, op \ang{0}\,N, \ang{90}\,E. (c) De resultaten verschillen: \omterm{def:b1:plate-kinematics:finite-rotation}{eindige rotaties} commuteren niet. Opeenvolgende reconstructies moeten worden samengesteld door de rotatiematrices in de juiste volgorde te vermenigvuldigen, nooit door hoeken of polen op te tellen.
\end{solution}

\begin{solution}{exo:b1:plate-kinematics:9}
(a) In de snelheidsruimte legt men $a$ in de oorsprong, $b$ \qty{30}{mm.yr^{-1}} naar het oosten en $c$ \qty{40}{mm.yr^{-1}} naar het zuiden. (b) $\vec v_{CB} =
\vec v_{CA} - \vec v_{BA} = (0, -40) - (30, 0) = (-30, -40)$ (oost- en noordcomponenten): \qty{50}{mm.yr^{-1}} naar het \omterm{def:b1:spherical-astronomy:horizontal}{azimut}
$\arctan\bigl((-30)/(-40)\bigr) + \ang{180} = \ang{216.9}$, zuidwest. (c) De rug staat loodrecht op deze richting: zij loopt
$\ang{216.9} - \ang{90} = \ang{126.9}$, dat is noordwest--zuidoost.
\end{solution}

\begin{solution}{exo:b1:plate-kinematics:10}
(a) In de matrix van Rodrigues draagt de eenheidsmatrix $3\cos\alpha$ bij, de matrix
$\vec n\,\vec n^{\mathsf T}$ het spoor $n_x^2 + n_y^2 + n_z^2 = 1$, maal
$1 - \cos\alpha$, en de antisymmetrische matrix niets: $\operatorname{tr}R =
3\cos\alpha + 1 - \cos\alpha = 1 + 2\cos\alpha$. (b) Hier is $\operatorname{tr}R = 0$,
dus $\cos\alpha = -1/2$ en $\alpha = \ang{120}$. De as voldoet aan $R\vec n =
\vec n$: $R(x, y, z) = (z, x, y)$, dus $x = y = z$ en $\vec n =
(1, 1, 1)/\sqrt3$, op \omterm{def:b1:spherical-astronomy:geographic}{breedtegraad} $\arcsin(1/\sqrt3) = \ang{35.26}$\,N en
\omterm{def:b1:spherical-astronomy:geographic}{lengtegraad} \ang{45}\,E. (c) $R\vec e_x = \vec e_y$, $R\vec e_y = \vec e_z$,
$R\vec e_z = \vec e_x$: de drie assen worden verwisseld, zoals een derde van een omwenteling om de diagonaal doet. De formule van Rodrigues met $\alpha = +\ang{120}$ om $\vec n$ geeft inderdaad $\vec e_x\mapsto -\tfrac12\vec e_x + \tfrac12(1, 1, 1) +
\tfrac{\sqrt3}{2}\cdot\tfrac{1}{\sqrt3}(0, 1, -1) = (0, 1, 0)$, dus de rotatie is tegen de klok in gezien van \ang{35.26}\,N, \ang{45}\,E.
\end{solution}

\begin{solution}{exo:b1:plate-kinematics:11}
(a) Op een transform staat de relatieve snelheid $\vec\omega\times\vec r$ langs de breuk, dus evenwijdig aan $\vec t$; zij staat loodrecht op $\vec\omega$, dus $\vec t\cdot\vec\omega = 0$. Het vlak loodrecht op $\vec t$ door het middelpunt bevat $\vec r$ (omdat $\vec t$ horizontaal is) en $\vec\omega$: het is het vlak van de grote cirkel door het punt, loodrecht op de breuk, en de pool ligt op die cirkel. Twee zulke vlakken snijden elkaar langs de lijn van $\vec t_1\times\vec t_2$. (b) $\vec t_1 = (0.5436, 0.8367, -0.0673)$ bij het knooppunt en $\vec t_2 = (0.4942, 0.8639, -0.0968)$ op IJsland, dus
$\vec t_1\times\vec t_2 = (-0.02285, 0.01936, 0.05612)$, met lengte $0.06362$.
Haar \omterm{def:b1:spherical-astronomy:geographic}{breedtegraad} is $\arcsin(0.05612/0.06362) = \ang{61.9}$\,N en haar \omterm{def:b1:spherical-astronomy:geographic}{lengtegraad} $\ang{180} - \arctan(0.01936/0.02285) = \ang{139.7}$\,E (negatieve $x$-component): de pool van het hoofdstuk. (c) Slecht: de twee grote cirkels kruisen elkaar onder een kleine hoek, want de richtingen verschillen maar \ang{8} over \ang{26} van de grens. Een verandering van één van beide richtingen met één graad verplaatst de pool met ongeveer \ang{15}. Werkelijke bepalingen combineren veel transforms met spreidingssnelheden, waarvan de variatie langs de grens de hoekafstand tot de pool vastlegt.
\end{solution}

\begin{solution}{exo:b1:plate-kinematics:12}
(a) $(94 - 4)/3 = \qty{30}{mm.yr^{-1}}$. (b) De verschillen over precies één jaar zijn
$24 - 4 = 20$, $25 - 5 = 20$, $26 - 6 = 20$, $65 - 15 = 50$, $74 - 24 = 50$,
$75 - 25 = 50$, $76 - 26 = 50$, $85 - 65 = 20$ en $94 - 74 = 20$ (mm). Gesorteerd:
vijf waarden van 20 en vier van 50; de mediaan is \qty{20}{mm.yr^{-1}}. (c) Twee posities die één jaar uit elkaar liggen zijn in dezelfde fase van de seizoensschommeling, die in hun verschil wegvalt. Een sprong raakt alleen de paren die er tussen liggen, hier vier van de negen; zolang minder dan de helft van de paren een sprong omvat, negeert de mediaan hem, terwijl de eindpuntschatting in (a) de hele sprong als beweging neemt en \qty{50}{\percent} te hoog uitkomt. De methode faalt wanneer sprongen zo vaak voorkomen dat de meeste paren van één jaar er één bevatten, of wanneer het station onregelmatig beweegt, bijvoorbeeld na een grote aardbeving.
\end{solution}

\begin{solution}{pb:b1:plate-kinematics:1}
\textbf{1.} $\omega = 0.210\times(\pi/180)/\num{e6} = \qty{3.665e-9}{rad.yr^{-1}}$;
$\omega R = \num{3.665e-9}\times\qty{6.371e9}{mm} = \qty{23.35}{mm.yr^{-1}}$.

\textbf{2.} $\cos\Delta = \sin\ang{61.8}\sin\ang{64.5} + \cos\ang{61.8}
\cos\ang{64.5}\cos(\ang{-157.6}) = 0.79546 - 0.18809 = 0.60737$, dus
$\Delta = \ang{52.60}$.

\textbf{3.} $v = 23.35\sin\ang{52.60} = 23.35\times0.7944 =
\qty{18.55}{mm.yr^{-1}}$.

\textbf{4.} De teller is $\sin\ang{157.6}\cos\ang{61.8} = 0.18011$ en de noemer
\[
  \cos\ang{64.5}\sin\ang{61.8} - \sin\ang{64.5}\cos\ang{61.8}\cos\ang{157.6} =
  0.37940 + 0.39436 = 0.77376;
\]
beide zijn positief, dus $\beta = \arctan 0.2328 =
\ang{13.1}$. Eurazië beweegt naar $\ang{13.1} + \ang{90} = \ang{103.1}$.

\textbf{5.} Oost: $18.55\sin\ang{103.1} = \qty{18.07}{mm.yr^{-1}}$; noord:
$18.55\cos\ang{103.1} = \qty{-4.20}{mm.yr^{-1}}$, dat is \qty{4.20}{mm.yr^{-1}}
naar het zuiden.

\textbf{6.} $\cos\Delta = 0.8851$, $\Delta = \ang{27.74}$, $v = 23.35
\sin\ang{27.74} = \qty{10.9}{mm.yr^{-1}}$. De twee platen draaien om een pool in oostelijk Siberië: de Gakkelrug ligt er dicht bij, en de snelheid daalt met $\sin\Delta$, hoewel het dezelfde platen zijn.

\textbf{7.} West: $94/9.779 = \qty{9.61}{mm.yr^{-1}}$; oost: $97/9.779 =
\qty{9.92}{mm.yr^{-1}}$.

\textbf{8.} $(94 + 97)/9.779 = 191/9.779 = \qty{19.53}{mm.yr^{-1}}$.

\textbf{9.} Uit C2An: $54.1/2.581 = \qty{20.96}{mm.yr^{-1}}$; uit de basis van
C1n: $19.8/0.781 = \qty{25.35}{mm.yr^{-1}}$. Zij komen niet overeen: hoe jonger de grens, hoe hoger de snelheid. Een echte verandering van snelheid is mogelijk, maar een stationair overschot in de afstanden, op alle leeftijden gelijk, zou precies dit doen.

\textbf{10.} De verhouding van de afstanden west tot oost is 0.96 tot 0.98: de oostelijke flank is een paar procent meer gegroeid. De volle snelheid hangt alleen af van de som van de twee afstanden, dus de asymmetrie verandert haar niet.

\textbf{11.} $(191 - 5.0)/9.779 = \qty{19.02}{mm.yr^{-1}}$;
$(54.1 - 5.0)/2.581 = \qty{19.02}{mm.yr^{-1}}$; $(19.8 - 5.0)/0.781 =
\qty{18.95}{mm.yr^{-1}}$. De drie stemmen nu overeen binnen \qty{0.1}{mm.yr^{-1}}. Dezelfde \qty{5}{km} is een kwart van de \qty{19.8}{km} van het jongste paar, maar minder dan \qty{3}{\percent} van de \qty{191}{km} van het oudste; oude grenzen geven de betrouwbaarder snelheden.

\textbf{12.} De gecorrigeerde halve snelheid is $186/9.779/2 = \qty{9.51}{mm.yr^{-1}}$,
en $9.51\times0.084 = \qty{0.80}{km}$, een derde van de \qty{2.5}{km} waarmee elke grens is verschoven. Men zou verwachten dat zo'n smalle streep uitgesmeerd is. Dat zij herkend wordt op de Reykjanes-rug, terwijl zij op de Gakkelrug ontbreekt, die uitgehongerd is naar magma, suggereert dat een ruime toevoer van magma het veld zelfs in smalle strepen trouw vastlegt.

\textbf{13.} De gecorrigeerde snelheden, 18.95 tot \qty{19.02}{mm.yr^{-1}}, stemmen overeen met de \qty{19.0}{mm.yr^{-1}} van de plaatvector; zonder de correctie overtroffen zij die met 3 tot 33 procent.

\textbf{14.} $\cos\Delta = \sin\ang{64.139}\sin\ang{59.018} +
\cos\ang{64.139}\cos\ang{59.018}\cos\ang{27.554} = 0.97054$, dus $\Delta =
\ang{13.94}$; de verbinding is $2R\sin(\Delta/2) = 12742\times\sin\ang{6.97} =
\qty{1546}{km}$.

\textbf{15.} $1000/17.5 = 57$ jaar.

\textbf{16.} De twee richtingen verschillen $\ang{123.1} - \ang{114.8} =
\ang{8.3}$: de horizontale component langs de grote cirkel is
$18.9\cos\ang{8.3} = \qty{18.70}{mm.yr^{-1}}$, en langs de verbinding
$18.70\cos\ang{6.97} = \qty{18.56}{mm.yr^{-1}}$, ongeveer \qty{18.6}{mm.yr^{-1}}.

\textbf{17.} De verbinding wordt \qty{17.5}{mm.yr^{-1}} langer, \qty{6}{\percent} minder dan de voorspelde \qty{18.6}{mm.yr^{-1}}. Beide stations liggen ver van de plaatgrens, dus het verschil is niet lokaal. Ofwel bewegen de platen vandaag iets trager dan gemiddeld over de laatste drie miljoen jaar, ofwel draagt één van de twee schattingen een kleine systematische fout; de metingen alleen kunnen dat niet uitmaken. Een \omterm{def:b1:plate-kinematics:euler-pole}{Eulervector} die alleen uit satellietpositionering is berekend, voorspelt \qty{17.1}{mm.yr^{-1}} voor deze verbinding, dicht bij de gemeten waarde.

\textbf{18.} Oost: $21.12 - (-1.07) = \qty{22.19}{mm.yr^{-1}}$; noord:
$-6.97 - (-0.97) = \qty{-6.00}{mm.yr^{-1}}$. Snelheid $\sqrt{22.19^2 + 6.00^2} =
\qty{23.0}{mm.yr^{-1}}$, naar $\ang{180} - \arctan(22.19/6.00) =
\ang{105.1}$.

\textbf{19.} De richting ligt dicht bij de \ang{103} van de plaatvector, maar de snelheid overtreft diens \qty{18.6}{mm.yr^{-1}} met een kwart. Beide stations liggen binnen de zone waar de grens vervormt, en Höfn, nabij de grootste ijskap, beweegt mee met het herstel van de \omterm{def:b1:layered-earth:layers}{korst} terwijl het ijs dunner wordt: omhoog met \qty{12.9}{mm.yr^{-1}}, en zijwaarts, weg van het ijs. Stations binnen de grenszone meten de plaatbeweging plus de lokale vervorming.

\textbf{20.} $0.210\times10 = \ang{2.10}$.

\textbf{21.} Het punt beweegt op een \omterm{def:b1:spherical-astronomy:spherical-triangle}{kleine cirkel} met straal $R\sin\Delta$:
$6371\times0.7944\times(2.10\times\pi/180) = \qty{185.5}{km}$. De gecorrigeerde breedte van de \omterm{def:b1:layered-earth:layers}{korst} die is gevormd sinds de jonge rand van C5n, \qty{186}{km} in \qty{9.8}{Myr}, is bijna dezelfde: de beweging is over tien miljoen jaar stationair geweest.

\textbf{22.} Componenten (oost, noord): Eurazië
$(22.9\sin\ang{95.1}, 22.9\cos\ang{95.1}) = (22.81, -2.04)$, Nubië
$(19.2\sin\ang{103.8}, 19.2\cos\ang{103.8}) = (18.65, -4.58)$. Eurazië ten opzichte
van Nubië: $(4.16, 2.54)$, dat is \qty{4.9}{mm.yr^{-1}} naar
$\arctan(4.16/2.54) = \ang{58.6}$, oost-noordoost.

\textbf{23.} Een knooppunt van drie riggen is stabiel, wat de snelheden ook zijn. De snelheid van vraag 22 laat Eurazië en Nubië uit elkaar bewegen, van oost-noordoost tot west-zuidwest: hun grens bij het knooppunt moet openen, langzaam; het is de Terceira-rift, wiens oost-zuidoost-richting de opening schuin maakt.

\textbf{24.} Deel I: \qty{18.6}{mm.yr^{-1}} in het midden van IJsland (\qty{19.0}{mm.yr^{-1}} bij het profiel van deel II). Deel II: 18.95 tot \qty{19.02}{mm.yr^{-1}} na de correctie. Deel III: de verbinding Reykjavik–Stavanger wordt \qty{17.5}{mm.yr^{-1}} langer, waar de plaatvector \qty{18.6}{mm.yr^{-1}} voorspelt; het paar binnen IJsland geeft \qty{23.0}{mm.yr^{-1}}, wat niet de plaatbeweging is. De plaatvector en de strepen komen overeen; de satellieten wijken \qty{6}{\percent} af. De strepen middelen over tot tien miljoen jaar, de plaatvector over \qty{3.16}{Myr}; de satellieten meten de beweging van de laatste twintig tot dertig jaar.

\textbf{25.} De plaatvector geeft \qty{18.55}{mm.yr^{-1}} over IJsland, en de gecorrigeerde strepen \qty{19.0}{mm.yr^{-1}} ernaast: op het dichtstbijzijnde millimeter per jaar is de \omterm{def:b1:plate-kinematics:spreading-rate}{volle spreidingssnelheid} over IJsland \textbf{\qty[text-series-to-math=true]{19}{mm.yr^{-1}}}.
\end{solution}
